% !TeX root = ../../main.tex
\section{\bilingual{Asymptotic Speeds of the Extremal Positions in \texorpdfstring{$N$}{N}-BRW of I.I.D. Bi-Branching Law}{同分布二分支\texorpdfstring{$N$}{N}-分支随机游走最值的渐近速率}}
    \bilingual{
        B\'erard and Gou\'er\'e \mycite{bérard2010brunetderridabehaviorbranchingselectionparticle} used Kingman's subadditive ergodic theorem to show that, for the binary branching law $\Xi=(p,p)$, if the branching law satisfies the condition that there exists $\lambda>0$ such that $\psi(\lambda)>0$ and $\psi(-\lambda)>0$, then the asymptotic speeds of the leftmost and rightmost positions of the $N$-branching random walk starting with $N$ particles exist in the sense of $\mathrm L^1$. This paper will follow this method, but under a more general branching law that can be infinitely branching, it will be shown that the asymptotic speed of the single-sided leftmost position of the $N$-branching random walk starting with at most $N$ particles exists in the sense of $\mathrm L^1$, see \Cref{corollary:leftmost-speed:N0}.

        Moreover, they also showed that this asymptotic speed has an order of $v_\infty(p)-\chi(p)\log^{-2}N+o(\log^{-2}N)$, where
        \begin{equation*}
            \chi(p)=\frac{\pi^2}{2}\lambda^*\psi''(\lambda^*),
        \end{equation*}
        and $\lambda^*>0$ satisfies $\lambda^*\psi'(-\lambda^*)-\psi(-\lambda^*)=\log 2$. The final result about the order of the asymptotic speed of the single-sided leftmost position of $N$-branching random walk will also be given in this thesis, see \Cref{theorem:speed}, but mainly using the conclusion of \nameref{theorem:mogulskii} and the method of A\"id\'ekon and Hu \mycite{aïdékon2025averageweightpercolationcompletegraph}.
    }{
        B\'erard 与 Gou\'er\'e \mycite{bérard2010brunetderridabehaviorbranchingselectionparticle} 利用 Kingman 次可加遍历定理得出, 在二分支的分支律$\Xi=(p,p)$下, 若分支律满足存在$\lambda>0$使得$\psi(\lambda)>0$与$\psi(-\lambda)>0$的条件, 则从$N$个粒子出发的$N$-分支随机游走最左侧位置与最右侧位置渐近速率在$\mathrm L^1$的意义下存在. 本文将遵照这一方法, 但在更一般的可无限分支的分支律下, 证明从少于等于$N$个粒子出发的$N$-分支随机游走单侧最左侧位置渐近速率在$\mathrm L^1$的意义下存在, 详见\Cref{corollary:leftmost-speed:N0}.

        另外, 他们还证得该渐近速率具有阶数$v_\infty(p)-\chi(p)\log^{-2}N+o(\log^{-2}N)$, 其中
        \begin{equation*}
            \chi(p)=\frac{\pi^2}{2}\lambda^*\psi''(\lambda^*),
        \end{equation*}
        其中$\lambda^*>0$满足$\lambda^*\psi'(-\lambda^*)-\psi(-\lambda^*)=\log 2$. 本文最后也将给出一个关于从少于等于$N$个粒子出发的$N$-分支随机游走单侧最左侧位置渐近速率阶数的结论, 详见\Cref{theorem:speed}, 但主要用到\nameref{theorem:mogulskii}的结论和A\"id\'ekon 与 Hu \mycite{aïdékon2025averageweightpercolationcompletegraph}的方法.
    }

\raggedbottom
